%% kapa: exam=2 round=22 year=2011 period=1 subject=appraisal_practice subject_ko=감정평가실무 seq=4 number=4 points=10
%% source: Q-Net 감정평가사 기출문제 (cst00309) artlSeq=5040389 file=extracted/1008/5040389/2011년 제22회 감정평가사 2차시험 기출문제/1교시.hwp
%% conversion: kordoc 4.24.51 -> tools/gen2.py (auto); review=proofread-sonnet; flags=tikz
\begin{longquestion}{【문제 4】}
서울지방법원 민사 00단독 재판장 판사 한공정은 민사소송사건의 심리를 위하여 다음과 같은 사건의 감정평가를 당신에게 의뢰하였다. 주어진 자료를 검토한 후 판사의 감정요청사항에 대하여 견해를 간단하게 약술하시오. (10점) \par % [원문 그대로] 00단독
<사건의 개요> \par
1) 감정평가 목적물 : 서울특별시 OO구 OO동 1669-1 한라산 오리엔탈 1층 13호 건물 29.00m\textsuperscript{2}, 대지권 5,500m\textsuperscript{2} × 5,500분지10.60m\textsuperscript{2} \par
2) 청구내용 : 원고 이대리는 2009. 09. 20 피고 (주)한라산개발이 시공하여 분양하는 위 목적물을 분양받아 현재 점포를 운영 중에 있으며, 피고는 분양 당시 전ㆍ후면 모두 인근 도로와 같은 높이의 평탄한 건축물인 것처럼 광고하였고 이에 원고는 그 중 후면 상가인 위 목적물을 분양 받았다. 그러나 건축물이 완공되어 입주하여 보니 건물의 전면 상가와는 달리 후면의 상가는 인근도로에 비해 약 1.2m 정도 높은 상태로 계단이 설치되어 고객의 통행이 불편한 구조로 되어 있음을 알게 되었고 분양 당시 광고와는 다른 구조로 인하여 실제 영업이익이 기대한 바에 미치지 못할 뿐만 아니라 현재 시점에서 위 목적물의 시가가 분양가에 미치지 못하는 등 손해를 입게 되었음을 주장하면서 이러한 손해를 이유로 2009. 09. 20 당시 분양가격이 과다하고 따라서 피고가 부당하게 얻은 분양가격의 일부를 반환할 것을 청구하는 소를 제기하였다. \par
물음(1) 원고의 주장이 타당한지 여부와 분양당시 전면 상가의 분양가격과 비교할 때 후면 상가의 분양가격이 적정한지 여부를 약술하시오. (5점) \par
물음(2) 분양가격이 적정하였다고 판단한다면 그 근거를 약술하고, 적정하지 않았다고 판단한다면 적정가격 수준은 어느 정도인지를 약술하시오. (5점) \par
<자료 1> 분양가격 자료(2009. 09. 20 당시 1층 상가의 분양자료) \par
{\small\centering
\begin{tabular}{|c|c|c|c|c|c|}
\hline
\multirow{2}{*}{\makecell{구분}} & \multicolumn{2}{|c|}{\makecell{면적(m\textsuperscript{2})}} & \multicolumn{2}{|c|}{\makecell{분양가격}} & \multirow{2}{*}{\makecell{비고}} \\
\cline{2-5}
 & \makecell{전용면적} & \makecell{대지권} & \makecell{금액(원)} & \makecell{단가(원/m\textsuperscript{2})} &  \\
\hline
\makecell{1층 4호} & \makecell{26.00} & \makecell{9.50} & \makecell{498,000,000} & \makecell{19,154,000} & \makecell{전면} \\
\hline
\makecell{1층 6호} & \makecell{65.00} & \makecell{23.80} & \makecell{1,200,000,000} & \makecell{18,462,000} & \makecell{전면} \\
\hline
\makecell{1층 9호} & \makecell{26.00} & \makecell{9.50} & \makecell{475,000,000} & \makecell{18,269,000} & \makecell{전면} \\
\hline
\makecell{1층 13호} & \makecell{29.00} & \makecell{10.60} & \makecell{425,000,000} & \makecell{14,655,000} & \makecell{후면} \\
\hline
\makecell{1층 14호} & \makecell{29.00} & \makecell{10.60} & \makecell{425,000,000} & \makecell{14,655,000} & \makecell{후면} \\
\hline
\makecell{1층 17호} & \makecell{32.00} & \makecell{11.70} & \makecell{460,000,000} & \makecell{14,375,000} & \makecell{후면} \\
\hline
\end{tabular}\par}

<자료 2> 실거래가격 자료 및 임대사례 \par
{\small\centering
\begin{tabular}{|c|c|c|c|c|c|c|}
\hline
\multirow{2}{*}{\makecell{구분}} & \multicolumn{3}{|c|}{\makecell{거래사례}} & \multicolumn{2}{|c|}{\makecell{임대료(천원)}} & \multirow{2}{*}{\makecell{임대\newline{}수익율}} \\
\cline{2-6}
 & \makecell{금액(천원)} & \makecell{매매일자} & \makecell{상승률} & \makecell{보증금} & \makecell{월임료} &  \\
\hline
\makecell{1층 4호} & \makecell{550,000} & \makecell{2011.08.04} & \makecell{10.44\%} & \makecell{-} & \makecell{-} & \makecell{자가사용} \\
\hline
\makecell{1층 6호} & \makecell{-} & \makecell{-} & \makecell{-} & \makecell{200,000} & \makecell{2,800} & \makecell{3.97\%} \\
\hline
\makecell{1층 9호} & \makecell{-} & \makecell{-} & \makecell{-} & \makecell{55,000} & \makecell{1,150} & \makecell{3.72\%} \\
\hline
\makecell{1층 13호} & \makecell{-} & \makecell{-} & \makecell{-} & \makecell{-} & \makecell{-} & \makecell{자가사용} \\
\hline
\makecell{1층 14호} & \makecell{460,000} & \makecell{2010.06.15} & \makecell{8.24\%} & \makecell{60,000} & \makecell{1,000} & \makecell{3.81\%} \\
\hline
\makecell{1층 17호} & \makecell{500,000} & \makecell{2010.10.31} & \makecell{8.70\%} & \makecell{70,000} & \makecell{1,050} & \makecell{3.80\%} \\
\hline
\end{tabular}\par}

※ 임대사례는 모두 2009년에 임대되었고 2011년 9월에 갱신되었음 \par
<자료 3> 경제동향 \par
{\small\centering
\begin{tabular}{|c|c|c|c|}
\hline
\makecell{구분} & \makecell{지가변동률} & \makecell{주택가격\newline{}상승률} & \makecell{생산자물가\newline{}상승률} \\
\hline
\makecell{2009.01.01 ～ 2009.12.31} & \makecell{1.07525} & \makecell{1.0473} & \makecell{1.0435} \\
\hline
\makecell{2010.01.01 ～ 2010.12.31} & \makecell{1.05366} & \makecell{1.0408} & \makecell{1.0459} \\
\hline
\makecell{2011.01.01 ～ 2011.08.31} & \makecell{1.03824} & \makecell{1.0311} & \makecell{1.0330} \\
\hline
\end{tabular}\par}

※ 지가변동률과 주택가격상승률은 OO구 평균임 \par
<자료 4> 상가배치도(축척 없음) \par
\begin{center}
\begin{tikzpicture}[x=0.0125cm,y=-0.0125cm,line width=0.6pt,font=\scriptsize\bfseries]
% roads and outer walls
\draw (66,42)--(66,452)--(22,452);
\draw (120,42)--(120,452)--(606,452)--(606,42);
\draw (660,42)--(660,452)--(708,452);
\draw (22,526)--(708,526);
% bottom shops
\draw (138,344) rectangle (579,432);
\foreach \x in {187,234,342,390,437,484,531} \draw (\x,344)--(\x,432);
\node at (163,390) {4}; \node at (210,390) {5}; \node at (288,390) {6};
\node at (413,390) {7}; \node at (460,390) {8}; \node at (507,390) {9}; \node at (555,390) {10};
\draw (342,400) arc[start angle=-90,end angle=0,x radius=24,y radius=32];
\draw (390,400) arc[start angle=-90,end angle=-180,x radius=24,y radius=32];
% left block 1-3
\draw (138,222)--(234,222)--(234,344);
\draw (138,222)--(138,344);
\draw (138,262)--(234,262); \draw (138,303)--(234,303);
\node at (187,242) {1}; \node at (187,283) {2}; \node at (187,323) {3};
% toilets, stairs, EV
\draw (254,222)--(463,222);
\draw (254,222)--(254,323)--(463,323);
\draw (254,270)--(330,270);
\draw (330,222)--(330,323);
\draw (383,222)--(383,323);
\node at (292,246) {여}; \node at (292,297) {남};
\draw (357,240)--(357,305);
\foreach \y in {243,263,283,303} \draw (345,\y)--(369,\y);
\draw (383,222)--(463,222);
\draw (393,225)--(408,258); \draw (423,225)--(408,258); \draw (423,225)--(438,258); \draw (453,225)--(438,258);
\draw (423,225)--(423,258); \draw (423,287)--(423,320);
\draw (393,320)--(408,287); \draw (423,320)--(408,287); \draw (423,320)--(438,287); \draw (453,320)--(438,287);
\node at (423,272) {E / V};
% right column
\draw (463,60) rectangle (579,101);
\node at (532,80) {17};
\draw (485,101)--(485,344); \draw (579,101)--(579,344); \draw (463,101)--(463,323);
\foreach \y/\n in {141/16,182/15,222/14,262/13,303/12,344/11} {}
\draw (485,141)--(579,141); \draw (485,182)--(579,182); \draw (485,222)--(579,222);
\draw (485,262)--(579,262); \draw (485,303)--(579,303); \draw (485,344)--(579,344);
\node at (532,121) {16}; \node at (532,161) {15}; \node at (532,202) {14};
\node at (532,242) {13}; \node at (532,283) {12}; \node at (532,324) {11};
% symbol at right
\draw (591,182)--(591,262); \draw (600,182)--(600,262); \draw (612,182)--(612,262); \draw (586,222)--(620,222);
\end{tikzpicture}
\end{center}
\end{longquestion}% [TikZ 재현] 상가배치도. E/V 칸 내부 사선은 원본 인쇄가 흐려 모양을 근사함 — 원본 대조 시 확인

